\section{本章小结}
\label{sec:topo-summary}

章首问的是：没有距离时，怎样谈论“附近”“趋于”和“连续”。本章的回答是把开集族作为唯一的基本数据。实轴开集满足的三条规则（空集与全空间开，任意并开，有限交开）被取作拓扑的定义；拓扑基用并生成全部开集，邻域与邻域基则只描述一个点周围的区域。离散拓扑、平凡拓扑和余可数拓扑表明，同一集合上的不同拓扑会给出截然不同的“附近”。

在固定的拓扑下，子集与环境的关系由闭集、内部、闭包与边界描述，它们都由邻域判据定义，并随拓扑改变；稠密集是闭包为全空间的集合，等价于与每个非空开集相交。点列收敛要求点列最终落入极限的每个邻域，收敛点列的子列收敛到同一点，闭集留住已有的极限；但余可数拓扑的例子说明，一般空间中点列不能反推闭包，因此闭包与稠密性都不以点列定义。

映射的连续性由开集的原像为开集来定义，等价于逐点的邻域判据与闭集原像判据，并保持点列极限，反过来的序列判据需要额外条件。开映射约束的是像，与连续性互不蕴含；同胚是双向连续的双射，连续双射为同胚当且仅当它是开映射。$(-1,1)$与$\mathbb R$同胚表明，拓扑不记录长度和有界性。子空间、有限乘积与商空间分别使包含映射、坐标投影与商映射连续，并给出映入子空间与乘积、映出商空间的连续性判据；$X\times[0,1]$是同伦的定义域，区间端点的粘合给出圆周。

空间的性质中，Hausdorff条件保证极限唯一，子空间与乘积保持它，商空间则可能破坏它，有两个原点的直线就是例子。第二可数是整体开集族的可数性，第一可数是逐点邻域族的可数性，可分是可数稠密点集的存在性，三者都不限制点的个数；第二可数蕴含另外两者，Sorgenfrey直线说明反向一般不成立。连通性只用开集表述，实轴上的连通集恰为区间，连续像保持连通并推出介值定理；路径连通蕴含连通，拓扑学家的正弦曲线说明反向不成立，局部路径连通时两种分支一致。紧性要求每个开覆盖都可有限化，连续像保持紧性，由此得到最值定理；从紧空间到Hausdorff空间的连续双射是同胚，这证明了区间粘合确实得到圆周。

本章所有概念都不依赖距离，这也划定了它的边界。距离会带来更多结构：闭包可用点到集合的距离刻画，第二可数与可分等价，紧性可用序列紧与“完备且全有界”判断，这些是第~\ref{chap:metric-spaces}章的内容；连续性的距离判据、均匀连续与映射列的收敛在第~\ref{chap:metric-space-maps}章讨论。另一方面，连续变形保留什么信息，要借助本章的乘积与商构造，在第~\ref{chap:homology}章用同调来回答。
